\documentclass[aspectratio=43,11pt]{beamer}
\usetheme{default}
\usecolortheme{default}
\definecolor{aggiemaroon}{RGB}{80,0,0}
\setbeamercolor{structure}{fg=aggiemaroon}
\setbeamertemplate{navigation symbols}{}
\setbeamertemplate{footline}[frame number]
\setbeamertemplate{itemize items}[circle]
\usepackage{amsmath,amssymb,booktabs}

\title{Decision and Risk Analysis}
\subtitle{\textbf{Practice Questions for Exam 2}}
\author{}
\date{
\color{aggiemaroon}{
DAEN 427 / ISEN 427 / ISEN 627 - Fall 2026\\
\textbf{Alexander Abuabara}
}
}
\titlegraphic{\vfill \includegraphics[height=1.2cm]{TAMU_logo.png}}

\newcommand{\divider}[2]{%
  \begin{frame}
    \centering
    {\Large #1}\\[0.6em]
    {\normalsize #2}
  \end{frame}}

\begin{document}

\begin{frame}
  \titlepage
\end{frame}

\begin{frame}{}
  \begin{itemize}
  \item Conceptual questions
  \item Open questions (for more understanding)
  \item Calculations (really short hand calculations)
  \item \textbf{Answers are at the end.}
  \end{itemize}
\end{frame}

%==================================================
\divider{Chapter 6}{Approximating the Posterior}

\begin{frame}{Question 6.1}
  For each MCMC diagnostic, say in plain words \textbf{what question it
  answers} and \textbf{what a good result looks like}.
  \vspace{0.6em}
  \begin{enumerate}
    \item Trace plot
    \item Density plots of the parallel chains, overlaid
    \item Effective sample size ratio $N_{\text{eff}}/N$
    \item $\hat R$
  \end{enumerate}
\end{frame}

\begin{frame}{Question 6.2}
  A grid of 101 values approximates the Beta(11, 3) posterior almost
  perfectly, while a 6-point grid only produced draws at 0.4, 0.6 and 0.8.
  \vspace{0.6em}

  \textbf{(a)} Why was the 6-point grid so poor? \\[0.3em]
  \textbf{(b)} If we keep 101 grid values \emph{per parameter}, how many
  posterior evaluations would a 3-parameter model need? A 10-parameter
  model? \\[0.3em]
  \textbf{(c)} What does this suggest about when to use MCMC instead?
\end{frame}

\begin{frame}{Question 6.3}
  Unlike grid approximation, MCMC draws are \emph{dependent}
  (each depends on the previous one) and are \emph{not} drawn directly
  from the posterior.
  \vspace{0.6em}

  \textbf{Why can we still use a Markov chain sample to estimate the
  posterior mean or a credible interval? What role does running the
  chain longer (and discarding burn-in) play?}
\end{frame}

\begin{frame}{Question 6.4}
  You run 4 parallel chains. Each trace plot, viewed \emph{alone}, looks
  like stable white noise. But chains 1--2 hover around $\theta \approx 0.3$
  and chains 3--4 hover around $\theta \approx 0.7$.
  \vspace{0.6em}

  \textbf{Would $\hat R$ be close to 1? Why do single-chain checks miss this
  problem, and what should you do?}
\end{frame}

\begin{frame}{Question 6.5}
  To ``reduce autocorrelation,'' a colleague keeps every 10th draw,
  cutting a fast-mixing \texttt{rstan} sample from 20{,}000 to 2{,}000 values.
  \vspace{0.6em}

  \textbf{Is this a good idea? What is being traded off, and when (if ever)
  is thinning worth it?}
\end{frame}

%==================================================
\begin{frame}{Question 6.6}
%  $Y \mid \pi \sim \text{Bin}(2, \pi)$, we observe $Y = 2$, and the prior is
%  Beta(2, 2), so $f(\pi) \propto \pi(1-\pi)$.
%  Use the grid $\pi \in \{0.25,\ 0.50,\ 0.75\}$.
%  \vspace{0.6em}
%
%  \textbf{(a)} At each grid value compute prior $\times$ likelihood
%  (up to a constant: $\pi(1-\pi) \cdot \pi^2$). \\[0.3em]
%  \textbf{(b)} Normalize to get the grid posterior probabilities. \\[0.3em]
%  \textbf{(c)} Compute the grid posterior mean and compare with the exact
%  posterior Beta(4, 2), mean $4/6$. Why do they differ?
%\end{frame}

  $Y \mid \pi \sim \text{Bin}(2, \pi)$, we observe $Y = 2$, and the prior is
  Beta(2, 2), so $f(\pi) \propto \pi(1-\pi)$.
  Use the grid $\pi \in \{0.25,\ 0.50,\ 0.75\}$.
  \vspace{0.6em}
 
  \textbf{(a)} At each grid value compute prior $\times$ likelihood.\\[0.3em]
  \textbf{(b)} Normalize to get the grid posterior probabilities. \\[0.3em]
  \textbf{(c)} Compute the grid posterior mean and compare with the exact
  posterior Beta(4, 2), mean $4/6$. Why do they differ?
  
  \vspace{.5cm}  
  \footnotesize You may drop constant factors (the 6 in the Beta(2, 2) density) that are the same for every $\pi$, so use $\pi(1-\pi) \cdot \pi^2$. % 
\end{frame}

\begin{frame}
\footnotesize 
\begin{align*}
\text{Prior: } & f(\pi) = 6\,\pi^{2-1}(1-\pi)^{2-1} \propto \pi(1-\pi) \\
\text{Likelihood: } & L(\pi \mid y=2) = \binom{2}{2}\pi^{2}(1-\pi)^{0} = \pi^{2} \\
\text{Posterior: } & f(\pi \mid y) \propto \pi(1-\pi)\cdot\pi^{2}
  = \pi^{3}(1-\pi) = \pi^{4-1}(1-\pi)^{2-1}\\
  & \;\Rightarrow\; \pi \mid y \sim \text{Beta}(4,2)
\end{align*}
\end{frame}

\begin{frame}{Question 6.7}
  You run 4 chains with 5{,}000 draws each (after burn-in).
  The variance of all draws combined is 0.09, and the typical
  variance within a single chain is 0.04. Compute
  $\hat R \approx \sqrt{\text{Var}_{\text{combined}} /
  \text{Var}_{\text{within}}}$ and interpret it.
\end{frame}

%==================================================
\divider{Chapter 7}{MCMC under the Hood}

\begin{frame}{Question 7.1}
  \textbf{(a)} Name the two steps of every Metropolis-Hastings iteration.
  \\[0.3em]
  \textbf{(b)} If the proposal is rejected, what is the chain's next
  value? Does that value count as a draw? \\[0.3em]
  \textbf{(c)} Why do these repeated values matter for getting the
  posterior right?
\end{frame}

\begin{frame}{Question 7.2}
  In the Normal-Normal example (posterior $\mu \mid y \sim N(4, 0.6^2)$),
  the Metropolis algorithm sometimes \emph{accepts} a proposal that is
  less plausible than the current value.
  \vspace{0.6em}

  \textbf{Why is that necessary? Describe the sample you would get from a
  rule that only accepts moves to more plausible values.}
\end{frame}

\begin{frame}{Question 7.3}
  The acceptance probability uses the \emph{unnormalized} posterior
  $f(\mu)L(\mu \mid y)$, not $f(\mu \mid y)$.
  \vspace{0.6em}

  \textbf{Why is this enough? And why is this the key reason MCMC works
  for models where grid approximation and exact formulas fail?}
\end{frame}

\begin{frame}{Question 7.4}
  Uniform proposal $\mu' \sim \text{Unif}(\mu - w,\ \mu + w)$.
  Compare $w = 0.01$, $w = 1$ and $w = 100$.
  \vspace{0.6em}

  \textbf{(a)} Which has the highest acceptance rate? \\[0.3em]
  \textbf{(b)} Sketch / describe the trace plot for $w = 0.01$ and
  $w = 100$. \\[0.3em]
  \textbf{(c)} Why is a very high acceptance rate \emph{not} a sign of a good
  chain?
\end{frame}

\begin{frame}{Question 7.5}
  For $\pi \in (0,1)$ we use an independence sampler: every proposal
  comes from the same Beta$(a,b)$ model, regardless of the current $\pi$.
  The acceptance probability is
  \[
    \alpha = \min\left\{1,\
      \frac{f(\pi' \mid y)}{f(\pi \mid y)} \cdot
      \frac{q(\pi)}{q(\pi')}\right\}.
  \]
  \textbf{(a)} Why is the correction $q(\pi)/q(\pi')$ needed if we propose
  from, say, Beta(5, 1)? \\[0.3em]
  \textbf{(b)} What happens to the correction with a Beta(1, 1) proposal?
\end{frame}

%\begin{frame}{Question 7.6}
%  Posterior: $\mu \mid y \sim N(4, 0.6^2)$, symmetric Uniform proposals
%  $\mu' \sim \text{Unif}(\mu - 2,\ \mu + 2)$, i.e.\ $w = 2$.
%  \vspace{0.6em}
%
%  \textbf{(a)} Compute $\alpha$ for a move from $\mu = 4$ to $\mu' = 5.2$,
%  and for the reverse move from $5.2$ to $4$. \\[0.3em]
%  \textbf{(b)} How does the ratio of these two move probabilities explain
%  why the chain ends up visiting values in proportion to the posterior?
%  
%\end{frame}
%
%\begin{frame}{Formulas for question 7.6}
%\footnotesize
%\begin{align*}
%\text{Posterior (kernel): } & f(\mu \mid y) \propto
%  \exp\!\left(-\frac{(\mu-4)^2}{2(0.6)^2}\right) \\[4pt]
%\text{Acceptance (symmetric proposal): } & \alpha =
%  \min\left\{1,\ \frac{f(\mu' \mid y)}{f(\mu \mid y)}\right\} \\[4pt]
%\text{Ratio for the Normal: } & \frac{f(\mu' \mid y)}{f(\mu \mid y)} =
%  \exp\!\left(-\frac{(\mu'-4)^2-(\mu-4)^2}{2(0.6)^2}\right) \\[4pt]
%\text{Decision: } & \text{draw } U \sim \text{Unif}(0,1);\\
%  & \text{next value} =
%  \begin{cases} \mu' & \text{if } U < \alpha \\ \mu & \text{otherwise} \end{cases}
%\end{align*}
%  
%\end{frame}

\begin{frame}{Question 7.6}
  The posterior is $\mu \mid y \sim N(0,\ 1)$: most
  plausible at $\mu = 0$. Metropolis proposes
  $\mu' \sim \text{Unif}(\mu - 3,\ \mu + 3)$ (symmetric), then accepts it
  with probability $\alpha$.
  \vspace{0.5em}

  \textbf{(a)} The chain is at $\mu = 0$ (the peak) and proposes
  $\mu' = 2$. Is 2 more or less plausible than 0? Compute $\alpha$.
  \\[0.3em]
  \textbf{(b)} The chain is at $\mu = 2$ and proposes $\mu' = 0$.
  Compute $\alpha$ \emph{without} a calculator. Why?
  \\[0.3em]
  \textbf{(c)} Out of 100 times the chain proposes these moves, about how
  many times will it actually move $0 \to 2$? And $2 \to 0$?
  What does that say about how often the chain visits 2 relative to 0?
  \vspace{0.3em}

  \footnotesize Formulas on the next slide.
\end{frame}

\begin{frame}{Formulas (for Question 7.6)}
  \footnotesize
  \begin{align*}
  \text{Posterior (kernel): } & f(\mu \mid y) \propto
    \exp\!\left(-\frac{\mu^2}{2}\right) \\[4pt]
  \text{Ratio: } &
    \frac{f(\mu' \mid y)}{f(\mu \mid y)} =
    \exp\!\left(-\frac{\mu'^2-\mu^2}{2}\right) \\[4pt]
  \text{Acceptance (symmetric): } & \alpha =
    \min\left\{1,\ \frac{f(\mu' \mid y)}{f(\mu \mid y)}\right\} \\[4pt]
  \text{Useful value: } & e^{-2} \approx 0.135
  \end{align*}
  \vspace{0.3em}
  \textbf{Rule of thumb:} moving to a \emph{more} plausible value is
  always accepted ($\alpha = 1$); moving to a \emph{less} plausible value
  is accepted with probability equal to the plausibility ratio.
\end{frame}

%==================================================
\begin{frame}{Question 7.7}
  Posterior: $\mu \mid y \sim N(4, 0.6^2)$, so
  $f(\mu \mid y) \propto \exp\!\left(-\dfrac{(\mu - 4)^2}{2(0.36)}\right)$.
  Symmetric Uniform proposals.
  \vspace{0.6em}

  The chain is at $\mu = 3.7$ and proposes $\mu' = 4.9$.
  \vspace{0.3em}

  \textbf{(a)} Compute the acceptance probability $\alpha$
  (hint: $e^{-1} \approx 0.368$). \\[0.3em]
  \textbf{(b)} To decide, we draw $U \sim \text{Unif}(0,1)$ and accept if
  $U < \alpha$. What is the next value if $U = 0.50$? If $U = 0.20$? \\[0.3em]
  \textbf{(c)} What would $\alpha$ be for the reverse move, $4.9 \to 3.7$?
\end{frame}

%\begin{frame}{Question 7.8}
%  Posterior Beta(3, 4): $f(\pi \mid y) \propto \pi^2 (1-\pi)^3$. \\
%  Proposal Beta(2, 1): $q(\pi) = 2\pi$ (proposes large $\pi$ more often).
%  \vspace{0.4em}
%
%  Current $\pi = 0.5$, proposal $\pi' = 0.25$.
%  \vspace{0.4em}
%
%  \textbf{(a)} Compute the posterior ratio $f(\pi' \mid y)/f(\pi \mid y)$.
%  \\[0.3em]
%  \textbf{(b)} Compute the correction $q(\pi)/q(\pi')$ and the acceptance
%  probability $\alpha$. \\[0.3em]
%  \textbf{(c)} The correction \emph{raised} $\alpha$. Explain why, in terms
%  of how often Beta(2, 1) proposes values like 0.25.
%\end{frame}

%==================================================
\divider{Chapter 8}{Posterior Inference \& Prediction}

\begin{frame}{Question 8.1}
  MoMA example: $\pi$ = proportion of artists who are Gen X or younger.
  Label each question as \textbf{estimation}, \textbf{hypothesis testing},
  or \textbf{prediction}:
  \vspace{0.6em}
  \begin{enumerate}
    \item What is a plausible range for $\pi$?
    \item Is $\pi$ below 0.20?
    \item How many of the next 20 artworks will be by Gen X artists?
    \item Is the typical value of $\pi$ closer to 0.15 or 0.30?
  \end{enumerate}
\end{frame}

\begin{frame}{Question 8.2}
  The middle 95\% credible interval for $\pi$ is $(0.10,\ 0.24)$.
  \vspace{0.6em}

  \textbf{(a)} Will the middle 99\% interval be wider or narrower? The
  middle 50\%? \\[0.3em]
  \textbf{(b)} Which credible level is the ``correct'' one to report?
\end{frame}

\begin{frame}{Question 8.3}
  MoMA example: prior Beta(4, 6), $Y = 14$ of $n = 100$ artists are Gen X
  or younger $\Rightarrow$ posterior Beta(18, 92). The middle 95\% credible
  interval for $\pi$ is $(0.10,\ 0.24)$.
  \vspace{0.6em}

  A student says: \emph{``If we repeated the study many times, 95\% of such
  intervals would contain $\pi$.''}
  \vspace{0.4em}

  \textbf{Is this the Bayesian interpretation? State it correctly. When
  would a highest posterior density interval differ from this one?}
\end{frame}

\begin{frame}{Question 8.4}
  Testing $H_a: \pi < 0.2$:
  \vspace{0.3em}
  \begin{center}
  \begin{tabular}{@{}lcc@{}}
    \toprule
    & Prior & Posterior \\
    \midrule
    $P(H_a)$ & 0.086 & 0.849 \\
    Odds & 0.094 & 5.62 \\
    \bottomrule
  \end{tabular}
  \end{center}
  Bayes factor $\approx 60$.
  \vspace{0.6em}

  \textbf{What does each number measure? Why can the Bayes factor be huge
  while the posterior probability is ``only'' 0.85?}
\end{frame}

\begin{frame}{Question 8.5}
  We want to test $H_0: \pi = 0.3$ vs.\ $H_a: \pi \neq 0.3$.
  \vspace{0.6em}

  \textbf{(a)} Why can't we compute $P(H_0 \mid y)$ or a Bayes factor
  directly? \\[0.3em]
  \textbf{(b)} What is the practical workaround, and what does it conclude
  here? \\[0.3em]
  \textbf{(c)} Why should that workaround be set \emph{before} seeing the
  data?
\end{frame}

\begin{frame}{Question 8.6}
  Predict $Y'$ = number of Gen X artists among 20 new works.
  A student uses the posterior mean $\hat\pi = 0.164$ and reports
  $Y' \sim \text{Bin}(20,\ 0.164)$.
  \vspace{0.6em}

  \textbf{(a)} Which source of variability is missing? \\[0.3em]
  \textbf{(b)} Will the student's predictions be too narrow or too wide? \\[0.3em]
  \textbf{(c)} With MCMC, why do we simulate one $Y'_i \sim
  \text{Bin}(20, \pi_i)$ per draw instead of using $20\,\pi_i$?
\end{frame}

\begin{frame}{Question 8.7}
  For $H_0: \pi \geq 0.2$ vs.\ $H_a: \pi < 0.2$:
  \begin{itemize}
    \item Frequentist p-value $= P(Y \le 14 \mid \pi = 0.2) \approx 0.08$
    \item Bayesian $P(H_a \mid Y = 14) \approx 0.849$
  \end{itemize}
  \vspace{0.4em}
  A student concludes: \emph{``There is an 8\% chance that $H_0$ is true.''}
  \vspace{0.6em}

  \textbf{What is wrong? What is the probability of $H_0$ according to the
  Bayesian analysis, and what does each method condition on?}
\end{frame}

%==================================================
\begin{frame}{Question 8.8}
  Prior $\pi \sim \text{Beta}(4, 6)$; data $Y = 14$ of $n = 100$.
  \vspace{0.3em}

  Recall: Beta$(\alpha, \beta)$ has mean $\dfrac{\alpha}{\alpha + \beta}$
  and mode $\dfrac{\alpha - 1}{\alpha + \beta - 2}$.
  \vspace{0.6em}

  \textbf{(a)} What are the posterior parameters? \\[0.3em]
  \textbf{(b)} Compute the posterior mean and mode. \\[0.3em]
  \textbf{(c)} Compute $E(Y' \mid Y = 14)$, the expected number of Gen X
  artists among 20 new works. Why can you get it without the full
  posterior predictive model?
\end{frame}

\begin{frame}{Question 8.9}
  Recall: odds $= \dfrac{p}{1 - p}$, \quad
  BF $= \dfrac{\text{posterior odds}}{\text{prior odds}}$.
  \vspace{0.6em}

  \textbf{(a)} A hypothesis has prior probability 0.2 and posterior
  probability 0.8. Compute the prior odds, posterior odds, and Bayes
  factor. \\[0.3em]
  \textbf{(b)} MoMA: prior odds of $H_a: \pi < 0.2$ are 0.094 and the
  Bayes factor is 60. Compute the posterior odds and the posterior
  probability $P(H_a \mid Y = 14)$.
\end{frame}

%==================================================
\divider{Chapter 9}{Simple Normal Regression}

\begin{frame}{Question 9.1}
  Bike model: posterior median slope $\beta_1 \approx 82.2$,
  80\% credible interval $(75.6,\ 88.8)$, and
  $P(\beta_1 > 0 \mid \text{data}) \approx 1$.
  \vspace{0.6em}

  \textbf{(a)} Interpret $\beta_1$ in one sentence. \\[0.3em]
  \textbf{(b)} Does this show that warmer weather \emph{causes} more
  ridership? What else could be going on?
\end{frame}

\begin{frame}{Question 9.2}
  Bike model: $Y_i \sim N(\beta_0 + \beta_1 X_i,\ \sigma^2)$,
  $X$ = temperature ($^\circ$F). The posterior median line is
  \[
    \hat\mu = -2194 + 82.2\,X .
  \]
  \textbf{(a)} Should we worry that the model predicts $-2194$ rides? \\[0.3em]
  \textbf{(b)} Why did we put the prior on the \emph{centered} intercept
  $\beta_{0c}$ (rides at an average temperature) instead of $\beta_0$?
\end{frame}

\begin{frame}{Question 9.3}
  From the posterior we get two 95\% intervals for 75$^\circ$F:
  \begin{itemize}
    \item Credible interval for the \emph{average} ridership on 75$^\circ$F
      days: $(3843,\ 4095)$
    \item Prediction interval for ridership on \emph{one} 75$^\circ$F day:
      $(1500,\ 6482)$
  \end{itemize}
  \vspace{0.4em}
  \textbf{(a)} Why is the second so much wider? \\[0.3em]
  \textbf{(b)} A manager must decide how many bikes to stock for
  tomorrow (75$^\circ$F). Which interval is relevant? \\[0.3em]
  \textbf{(c)} With 50{,}000 days of data, what would happen to each?
\end{frame}

\begin{frame}{Question 9.4}
  Data arrive in phases: 30 days, 60 days, then all 500 days.
  After the first 30 days (all cold winter days), the posterior for
  $\beta_1$ is still vague and does not reveal the warm-weather pattern.
  By 500 days it is concentrated near 80.
  \vspace{0.6em}

  \textbf{Why are the first 30 days so uninformative about the slope,
  even though 30 is not a tiny sample? What does this show about how a
  posterior evolves?}
\end{frame}

%==================================================
\begin{frame}{Question 9.5}
  \textbf{(a)} Prior $\beta_1 \sim N(100, 40^2)$. Using mean $\pm$ 2 SD,
  what range of slopes does the prior consider plausible? What does that
  mean in rides? \\[0.3em]
  \textbf{(b)} Prior $\sigma \sim \text{Exp}(0.0008)$. What is the prior
  mean of $\sigma$? (Exp$(\lambda)$ has mean $1/\lambda$.) \\[0.3em]
  \textbf{(c)} With the posterior median line $-2194 + 82.2X$: predict
  ridership at 75$^\circ$F, and the expected difference between a
  60$^\circ$F and a 70$^\circ$F day.
\end{frame}

\begin{frame}{Question 9.6}
  Three posterior draws, each with a standard Normal value $z$:
  \vspace{0.3em}
  \begin{center}
  \begin{tabular}{@{}ccccc@{}}
    \toprule
    Draw & $\beta_0$ & $\beta_1$ & $\sigma$ & $z$ \\
    \midrule
    1 & $-2200$ & 82 & 1300 & $\phantom{-}0.5$ \\
    2 & $-2100$ & 81 & 1250 & $-1.0$ \\
    3 & $-2300$ & 84 & 1280 & $\phantom{-}0.2$ \\
    \bottomrule
  \end{tabular}
  \end{center}
  \textbf{(a)} For each draw compute $\mu = \beta_0 + 75\beta_1$ and
  $Y' = \mu + \sigma z$. \\[0.3em]
  \textbf{(b)} Which set of values ($\mu$'s or $Y'$'s) builds the credible
  interval for the \emph{average} 75$^\circ$F day, and which builds the
  prediction interval? Which set is more spread out, and why?
\end{frame}

%==================================================
\divider{Chapter 15}{Hierarchical Models}

\begin{frame}{Question 15.1}
  Which data sets are grouped (hierarchical)? Name the group, and say
  why observations within a group are likely correlated.
  \vspace{0.6em}
  \begin{enumerate}
    \item Test scores of students from 12 schools
    \item Blood pressure measured monthly on 40 patients
    \item One survey response from each of 500 randomly sampled adults
    \item Daily high temperatures in 10 cities for a year
  \end{enumerate}
\end{frame}

\begin{frame}{Question 15.2}
  Cherry Blossom data: 252 race times from 36 runners (50s--60s), each
  observed in several years.
  \vspace{0.3em}

  Complete pooling (one regression of time on age): slope $\approx 0.27$
  min/year, 80\% CI $(-0.30,\ 0.84)$. Yet nearly \emph{every} runner gets
  slower with age, faster than this.
  \vspace{0.6em}

  \textbf{How can pooling all the data hide a trend that almost every
  individual shows? Which assumption is violated?}
\end{frame}

\begin{frame}{Question 15.3}
  A no-pooled model fits each runner separately. Runner 1 has only 5
  races and improved over them; the model predicts a time under 75
  minutes at age 62. The book considers 75--85 minutes more reasonable.
  \vspace{0.6em}

  \textbf{(a)} Why is the no-pooled prediction implausible? \\[0.3em]
  \textbf{(b)} Why can't a no-pooled model predict anything for a new
  runner (e.g.\ \emph{you})?
\end{frame}

\begin{frame}{Question 15.4}
  Runner 29 is fast and very consistent; runner 17 is slow and highly
  variable.
  \vspace{0.6em}

  \textbf{(a)} Distinguish \emph{within-group} and \emph{between-group}
  variability using these runners. \\[0.3em]
  \textbf{(b)} Why are a runner's own race times correlated? \\[0.3em]
  \textbf{(c)} How does a hierarchical (partial pooling) model use both
  kinds of variability to sit between the other two approaches?
\end{frame}

%==================================================
\begin{frame}{Question 15.5}
  Two runners, three races each:
  \vspace{0.3em}
  \begin{center}
  \begin{tabular}{@{}lccc@{}}
    \toprule
    & Ages & Times (min) & \\
    \midrule
    Runner A & 55, 56, 57 & 80, 81, 82 & \\
    Runner B & 60, 61, 62 & 70, 71, 72 & \\
    \bottomrule
  \end{tabular}
  \end{center}
  \textbf{(a)} No pooling: what is each runner's slope (min per year)? \\[0.3em]
  \textbf{(b)} Complete pooling: compute the least-squares slope through
  all 6 points,
  $b = \dfrac{\sum (x_i - \bar x)(y_i - \bar y)}{\sum (x_i - \bar x)^2}$.
  \\[0.3em]
  \textbf{(c)} Explain the contradiction.
\end{frame}

\begin{frame}{Question 15.6}
  A useful simplification: a hierarchical estimate for a runner is roughly
  \[
    w \times (\text{runner's own mean}) + (1 - w) \times (\text{overall mean}),
    \quad 0 \le w \le 1 .
  \]
  Overall mean: 90 min. Runners X and Y both have own mean 80 min;
  X ran 2 races, Y ran 10.
  \vspace{0.4em}

  \textbf{(a)} Compute the estimate with $w = 0.5$ and with $w = 0.9$.
  \\[0.3em]
  \textbf{(b)} Which $w$ fits X and which fits Y? Why? \\[0.3em]
  \textbf{(c)} Which values of $w$ give complete pooling and no pooling?
\end{frame}

%==================================================
\divider{Multiple Choice}{}

\begin{frame}{MCQ 1}
  You run 4 parallel chains. Which diagnostic most directly detects that
  the chains are exploring \emph{different} regions of the posterior?
  \vspace{0.6em}
  \begin{enumerate}[(A)]
    \item The effective sample size ratio $N_{\text{eff}}/N$
    \item The autocorrelation plot of chain 1
    \item $\hat R$
    \item The trace plot of a single chain
  \end{enumerate}
\end{frame}

\begin{frame}{MCQ 2}
  Metropolis algorithm with a symmetric proposal. The unnormalized
  posterior is $0.4$ at the current value and $0.1$ at the proposal.
  What is the probability of accepting the proposal?
  \vspace{0.6em}
  \begin{enumerate}[(A)]
    \item 0.10
    \item 0.25
    \item 0.40
    \item 1 --- it cannot be computed without normalizing
  \end{enumerate}
\end{frame}

\begin{frame}{MCQ 3}
  Which statement about the 95\% posterior prediction interval for
  ridership on a single 75$^\circ$F day is correct?
  \vspace{0.6em}
  \begin{enumerate}[(A)]
    \item It is narrower than the 95\% credible interval for the
      \emph{average} ridership on 75$^\circ$F days.
    \item It reflects only uncertainty in $\beta_0$ and $\beta_1$.
    \item It reflects both uncertainty in the parameters and day-to-day
      variability around the regression line.
    \item With enough data it shrinks to a single number.
  \end{enumerate}
\end{frame}

\begin{frame}{MCQ 4}
  Using the Cherry Blossom data, we want to predict race times for a
  \emph{new} runner who is not in the sample. Which models can do this?
  \vspace{0.6em}
  \begin{enumerate}[(A)]
    \item Only the no-pooled model
    \item Only the complete-pooled model
    \item The complete-pooled and hierarchical (partial pooling) models
    \item All three models
  \end{enumerate}
\end{frame}

%==================================================
\divider{Answers}{}

%---------------- Chapter 6 answers
\begin{frame}{Answer 6.1}
  \begin{enumerate}
    \item \textbf{Trace plot:} is the chain wandering stably? Good: looks
      like random noise around a stable level, with no trends or flat
      stretches.
    \item \textbf{Overlaid densities:} do the chains agree? Good: the
      density curves nearly coincide.
    \item \textbf{$N_{\text{eff}}/N$:} how many independent draws are our
      dependent draws worth? Good: not too small (above $\approx 0.1$).
    \item \textbf{$\hat R$:} is the variability across chains similar to the
      variability within chains? Good: close to 1 (below $\approx 1.05$).
  \end{enumerate}
\end{frame}

\begin{frame}{Answer 6.2}
  \textbf{(a)} The posterior can only put mass on grid points. With 6
  points, values between them (e.g.\ 0.7--0.9, where Beta(11, 3) is
  concentrated) can't be represented, so the approximation is crude.
  \vspace{0.4em}

  \textbf{(b)} $101^3 \approx 1.03$ million evaluations;
  $101^{10} \approx 1.1 \times 10^{20}$. Impossible.
  \vspace{0.4em}

  \textbf{(c)} The \emph{curse of dimensionality}: grids waste effort on
  implausible regions. MCMC spends its time where the posterior has mass,
  so it scales to many parameters.
\end{frame}

\begin{frame}{Answer 6.3}
  \begin{itemize}
    \item The chain is built so that, in the long run, the
      \emph{distribution of visited values} converges to the posterior.
    \item Averages over the chain (mean, quantiles) therefore approximate
      posterior quantities.
    \item Dependence doesn't make the estimates wrong; it makes each draw
      carry \emph{less information} (see effective sample size).
    \item Longer chains: more information and time to settle.
      Burn-in: discard early values that still reflect the arbitrary
      starting point.
  \end{itemize}
\end{frame}

\begin{frame}{Answer 6.4}
  \textbf{No! $\hat R$ would be far above 1} (well past 1.05).
  \begin{itemize}
    \item $\hat R \approx \sqrt{\text{Var}_{\text{combined}} /
      \text{Var}_{\text{within}}}$. The combined chains vary much more than
      any single chain.
    \item Each chain looks stable because it is stuck in its own region;
      only \emph{comparing} chains reveals they disagree.
    \item Action: don't trust the results. Run longer, use dispersed
      starting values, and check the model (e.g.\ a multimodal or poorly
      identified posterior).
  \end{itemize}
\end{frame}

\begin{frame}{Answer 6.5}
  \textbf{Usually not.}
  \begin{itemize}
    \item Thinning lowers autocorrelation but \emph{throws away}
      information: 20,000 correlated draws generally beat 2,000 less
      correlated ones.
    \item The chain is already fast-mixing, so there is little to gain.
    \item \texttt{rstan}/\texttt{rstanarm} (HMC) are efficient; their
      authors advise thinning only when memory/storage is a problem.
    \item To fix slow mixing, prefer running \emph{longer chains} (or
      improving the model).
  \end{itemize}
\end{frame}

\begin{frame}{Answer 6.6}
  \begin{center}
  \begin{tabular}{@{}cccc@{}}
    \toprule
    $\pi$ & $\pi^3(1-\pi)$ & $\times 256$ & Posterior \\
    \midrule
    0.25 & 0.0117 & 3 & $3/46 \approx 0.065$ \\
    0.50 & 0.0625 & 16 & $16/46 \approx 0.348$ \\
    0.75 & 0.1055 & 27 & $27/46 \approx 0.587$ \\
    \midrule
    & & 46 & 1 \\
    \bottomrule
  \end{tabular}
  \end{center}
  
%  \footnotesize
%  Why $256$? Any positive constant works: it multiplies every value
%  \emph{and} the total, so it cancels when we normalize (as does the 6 we
%  dropped). $256 = 4^4$ just turns the values into whole numbers.
%  \normalsize
  
  \textbf{(c)} Mean $= (0.25 \cdot 3 + 0.50 \cdot 16 + 0.75 \cdot 27)/46
  = 29/46 \approx 0.630$, vs.\ exact $0.667$.
  \vspace{0.3em}

  The grid can't put mass above 0.75, where Beta(4, 2) still has a lot
  of its mass (its mode is 0.75), so the coarse grid pulls the mean down.
  A finer grid fixes this.
\end{frame}

\begin{frame}{Answer 6.7}
  $\hat R \approx \sqrt{0.09/0.04} = \sqrt{2.25} = 1.5$.
  Far above 1 (or even a upper limit of 1.05): the chains together vary much more than each one
  alone, so they disagree. The simulation is unstable and its results
  should not be trusted.
\end{frame}

%---------------- Chapter 7 answers
\begin{frame}{Answer 7.1}
  \textbf{(a)} (1) \emph{Propose} a new value from a proposal model.
  (2) \emph{Decide}: accept it with probability $\alpha$, otherwise stay.
  \vspace{0.4em}

  \textbf{(b)} The chain stays put: the next value equals the current
  value, and yes, it counts as a draw in the sample.
  \vspace{0.4em}

  \textbf{(c)} Repeats are how the chain gives more weight to plausible
  values: it ``lingers'' where the posterior is high. Dropping the
  repeats would over-represent less plausible values and distort the
  approximation.
\end{frame}

\begin{frame}{Answer 7.2}
  \begin{itemize}
    \item The goal is not to \emph{find} the most plausible value but to
      visit every value \emph{in proportion to} its plausibility.
    \item Less plausible regions (the tails) must be visited sometimes,
      just less often. Accepting with probability equal to the
      plausibility ratio does exactly that.
    \item An ``uphill only'' rule acts like an optimizer: the chain climbs
      to $\mu \approx 4$ and stays there. The sample has almost no spread,
      so credible intervals would be far too narrow.
  \end{itemize}
\end{frame}

\begin{frame}{Answer 7.3}
  \begin{itemize}
    \item $\alpha$ depends only on a \emph{ratio}:
    \[
      \frac{f(\mu' \mid y)}{f(\mu \mid y)}
      = \frac{f(\mu')L(\mu' \mid y)/f(y)}{f(\mu)L(\mu \mid y)/f(y)}
      = \frac{f(\mu')L(\mu' \mid y)}{f(\mu)L(\mu \mid y)}.
    \]
      The normalizing constant $f(y)$ cancels.
    \item $f(y)$ is the hard part: a (multiple) integral over all
      parameters. Prior $\times$ likelihood is easy to evaluate at any
      point.
    \item So MCMC needs only what we can always compute, even in models
      with many parameters.
  \end{itemize}
\end{frame}

\begin{frame}{Answer 7.4}
  \textbf{(a)} $w = 0.01$: proposals are nearly identical to the current
  value, so almost all are accepted.
  \vspace{0.4em}

  \textbf{(b)} $w = 0.01$: a slowly wandering, smooth trace that explores
  only a small region (high autocorrelation). $w = 100$: long flat
  stretches where the chain is stuck, because most proposals land in
  implausible regions and are rejected.
  \vspace{0.4em}

  \textbf{(c)} Accepting tiny steps isn't exploring. Efficiency depends on
  \emph{how far} the chain moves, not just \emph{how often}. $w = 1$
  balances step size and acceptance (the ``Goldilocks'' choice).
\end{frame}

\begin{frame}{Answer 7.5}
  \textbf{(a)} Beta(5, 1) proposes large $\pi$ values much more often than
  small ones. Without correction, the chain would over-visit large values
  just because they are proposed more, biasing the sample toward the
  proposal. The factor $q(\pi)/q(\pi')$ penalizes moves to
  frequently proposed values and boosts moves to rarely proposed ones.
  \vspace{0.6em}

  \textbf{(b)} Beta(1, 1) is Uniform: $q(\pi) = q(\pi') = 1$, so the
  correction equals 1 and the algorithm reduces to the plain Metropolis
  ratio.
\end{frame}

%\begin{frame}{Answer 7.6}
%  \textbf{(a)} Proposal is symmetric, so $\alpha$ is the posterior ratio:
%  \[
%    \frac{f(5.2 \mid y)}{f(4 \mid y)}
%    = \exp\!\left(-\frac{(5.2-4)^2}{2(0.6)^2}\right) = e^{-2} \approx 0.135.
%  \]
%  $4 \to 5.2$: $\alpha \approx 0.135$. \quad
%  $5.2 \to 4$: more plausible, so $\alpha = 1$.
%  \vspace{0.4em}
%
%  \textbf{(b)} With $w = 2$ each value is a possible proposal from the
%  other, with equal proposal density, so
%  \[
%    \frac{P(4 \to 5.2)}{P(5.2 \to 4)} = 0.135 = \frac{f(5.2 \mid y)}{f(4 \mid y)}.
%  \]
%  The flow between any two values matches their relative posterior
%  plausibility, so in the long run the chain visits them in that
%  proportion.
%\end{frame}

\begin{frame}{Answer 7.6}
  \textbf{(a)} 2 is farther from the peak, so \emph{less} plausible:
  \[
    \frac{f(2 \mid y)}{f(0 \mid y)}
    = \exp\!\left(-\frac{2^2 - 0^2}{2}\right) = e^{-2} \approx 0.135
    \;\Rightarrow\; \alpha \approx 0.135
  \]
  \textbf{(b)} 0 is the peak, so it is \emph{more} plausible than 2:
  the ratio is $e^{2} > 1$, so $\alpha = 1$. No calculation needed:
  uphill moves are always accepted.
  \vspace{0.4em}

  \textbf{(c)} About $100 \times 0.135 \approx 13$ moves $0 \to 2$, and
  all $100$ moves $2 \to 0$.
  \vspace{0.2em}

  The chain moves from 2 back to 0 far more readily than from 0 out to 2.
  So in the long run it spends about $0.135$ times as much time at 2 as
  at 0, which is exactly $f(2 \mid y)/f(0 \mid y)$. The chain visits each
  value in proportion to its posterior plausibility.
\end{frame}

\begin{frame}{Answer 7.7}
  \textbf{(a)} Symmetric proposal, so $\alpha$ is the posterior ratio:
  \[
    \frac{f(4.9 \mid y)}{f(3.7 \mid y)}
    = \exp\!\left(-\frac{0.9^2 - 0.3^2}{2(0.36)}\right)
    = e^{-0.72/0.72} = e^{-1} \approx 0.368 .
  \]
  So $\alpha = \min\{1, 0.368\} = 0.368$.
  \vspace{0.4em}

  \textbf{(b)} $U = 0.50 > 0.368$: reject, next value is $3.7$.
  $U = 0.20 < 0.368$: accept, next value is $4.9$.
  \vspace{0.4em}

  \textbf{(c)} $3.7$ is more plausible than $4.9$ (closer to 4), so the
  ratio is $e^{1} > 1$ and $\alpha = 1$: always accepted.
\end{frame}

%\begin{frame}{Answer 7.8}
%  \textbf{(a)}
%  $\dfrac{0.25^2 \cdot 0.75^3}{0.5^2 \cdot 0.5^3}
%   = \dfrac{0.0625 \times 0.4219}{0.25 \times 0.125}
%   = \dfrac{0.02637}{0.03125} \approx 0.844$.
%  \vspace{0.4em}
%
%  \textbf{(b)} $q(0.5)/q(0.25) = 1.0/0.5 = 2$, so
%  $\alpha = \min\{1,\ 0.844 \times 2\} = \min\{1, 1.69\} = 1$.
%  Always accepted (without the correction it would be 0.844).
%  \vspace{0.4em}
%
%  \textbf{(c)} Beta(2, 1) rarely proposes small values like 0.25. To still
%  visit them as often as the posterior requires, the algorithm must accept
%  them more readily when they do come up. The correction does exactly
%  that.
%\end{frame}

%---------------- Chapter 8 answers
\begin{frame}{Answer 8.1}
  \begin{enumerate}
    \item Plausible range for $\pi$: \textbf{estimation}
      (credible interval).
    \item Is $\pi < 0.20$? \textbf{Hypothesis testing}
      (posterior probability, odds, Bayes factor).
    \item Next 20 artworks: \textbf{prediction}
      (posterior predictive model for $Y'$).
    \item Typical value of $\pi$: \textbf{estimation}
      (posterior mean, median or mode).
  \end{enumerate}
  \vspace{0.4em}
  Note: 1, 2 and 4 are about the parameter $\pi$; 3 is about future
  \emph{data}.
\end{frame}

\begin{frame}{Answer 8.2}
  \textbf{(a)} 99\%: \textbf{wider}, because it must hold more posterior
  probability: $(0.085,\ 0.265)$. \\
  50\%: \textbf{narrower}: $(0.139,\ 0.186)$.
  \vspace{0.6em}

  \textbf{(b)} None. 95\% is a convention, not a rule. The level should fit
  the purpose, and reporting several levels (or the whole posterior)
  communicates uncertainty better than any single interval.
\end{frame}
\begin{frame}{Answer 8.3}
  \textbf{Nope}. That is the frequentist confidence-interval
  interpretation (about the procedure over repeated samples).
  \vspace{0.4em}

  \textbf{Bayesian:} given the prior and the observed data, there is a
  95\% posterior probability that $\pi$ is between 0.10 and 0.24.
  The probability statement is about $\pi$, not about the interval.
  \vspace{0.4em}

  \textbf{HPD vs.\ middle:} they differ when the posterior is strongly
  \emph{skewed}. HPD keeps the most plausible values; the middle interval
  cuts equal tail areas. Beta(18, 92) is fairly symmetric, so they are
  nearly the same here.
\end{frame}

\begin{frame}{Answer 8.4}
  \begin{itemize}
    \item \textbf{Posterior probability 0.849:} our \emph{final} belief in
      $H_a$, combining prior and data.
    \item \textbf{Bayes factor}
      $= \dfrac{\text{posterior odds}}{\text{prior odds}}
      = \dfrac{5.62}{0.094} \approx 60$:
      how much the \emph{data} shifted the odds. A measure of evidence,
      not of belief.
    \item The prior was skeptical of $H_a$ (only 0.086). The data raised
      the odds 60-fold, but starting from a low base, the posterior ends
      at 0.85.
    \item Read both with nuance; avoid rigid ``significant'' cutoffs.
  \end{itemize}
\end{frame}

\begin{frame}{Answer 8.5}
  \textbf{(a)} $\pi$ is continuous, so $P(\pi = 0.3) = 0$ under both prior
  and posterior. Posterior odds and Bayes factor are undefined (0 / 0).
  \vspace{0.4em}

  \textbf{(b)} Define a practically meaningful buffer, e.g.\
  $H_0: \pi \in (0.25, 0.35)$. The whole buffer lies above the 95\% CI
  $(0.10, 0.24)$ $\Rightarrow$ strong evidence against $H_0$.
  \vspace{0.4em}

  \textbf{(c)} Choosing the buffer after seeing the data lets us move the
  goalposts to reach a desired conclusion. It should reflect subject-matter
  meaning, not results.
\end{frame}

\begin{frame}{Answer 8.6}
  \textbf{(a)} \emph{Posterior variability}: uncertainty in $\pi$ itself.
  Plugging in $\hat\pi$ keeps only sampling variability.
  \vspace{0.4em}

  \textbf{(b)} Too narrow (overconfident). Same center ($\approx 3.3$), but
  \begin{center}
  \begin{tabular}{@{}lc@{}}
    \toprule
    & Var$(Y')$ \\
    \midrule
    Bin(20, 0.164) & 2.74 \\
    Posterior predictive (Beta-Binomial) & 3.21 \\
    \bottomrule
  \end{tabular}
  \end{center}
  \textbf{(c)} $20\,\pi_i$ is an expected count and drops sampling
  variability. Drawing $Y'_i \sim \text{Bin}(20, \pi_i)$ for every
  posterior draw $\pi_i$ captures \emph{both} sources, approximating
  $f(y' \mid y) = \int f(y' \mid \pi) f(\pi \mid y)\, d\pi$.
\end{frame}

\begin{frame}{Answer 8.7}
  \textbf{The p-value is not $P(H_0 \mid \text{data})$.}
  \begin{itemize}
    \item p-value: \emph{assumes} $\pi = 0.2$ and asks how probable data
      at least this extreme are: $P(\text{data} \mid H_0)$.
    \item Bayesian: conditions on the observed data and asks how probable
      the hypothesis is: $P(H_0 \mid Y = 14) = 1 - 0.849 = 0.151$.
    \item These answer different questions and generally differ.
      The Bayesian answer is the one the student was trying to state.
  \end{itemize}
\end{frame}

\begin{frame}{Answer 8.8}
  \textbf{(a)} Beta$(4 + 14,\ 6 + 100 - 14) = $ Beta(18, 92).
  \vspace{0.4em}

  \textbf{(b)} Mean $= 18/110 \approx 0.164$; \quad
  mode $= 17/108 \approx 0.157$.
  \vspace{0.4em}

  \textbf{(c)} $E(Y' \mid Y = 14) = 20 \times E(\pi \mid Y = 14)
  = 20 \times 18/110 \approx 3.27$ (the book's 3.273).
  \vspace{0.3em}

  For any $\pi$, $E(Y' \mid \pi) = 20\pi$. Averaging over the posterior
  just replaces $\pi$ by its posterior mean. That shortcut works for the
  \emph{mean} only: probabilities and intervals for $Y'$ still need the
  full posterior predictive model (Q8.4).
\end{frame}

\begin{frame}{Answer 8.9}
  \textbf{(a)} Prior odds $= 0.2/0.8 = 0.25$; posterior odds
  $= 0.8/0.2 = 4$; BF $= 4/0.25 = 16$. The data made the hypothesis's
  odds 16 times larger.
  \vspace{0.6em}

  \textbf{(b)} Posterior odds $= 0.094 \times 60 \approx 5.64$.
  \[
    P(H_a \mid Y = 14) = \frac{\text{odds}}{1 + \text{odds}}
    = \frac{5.64}{6.64} \approx 0.849,
  \]
  matching the book.
\end{frame}

%---------------- Chapter 9 answers
\begin{frame}{Answer 9.1}
  \textbf{(a)} On average, each 1$^\circ$F increase in temperature is
  associated with about 82 more riders per day (80\% credible interval:
  about 76 to 89).
  \vspace{0.6em}

  \textbf{(b)} No. $P(\beta_1 > 0) \approx 1$ says the
  \emph{association} is almost surely positive, not that it is causal. The
  data are observational: warm days also tend to be in seasons with
  longer daylight, tourism, holidays and school schedules, any of which
  could affect ridership. Strong evidence of a relationship is not
  evidence of cause.
\end{frame}

\begin{frame}{Answer 9.2}
  \textbf{(a)} No. $-2194$ is the predicted ridership at 0$^\circ$F, far
  outside the observed temperatures in D.C. It is an extrapolation, so
  $\beta_0$ has no practical meaning. The line is only trusted where we
  have data.
  \vspace{0.4em}

  \textbf{(b)} We can reason about $\beta_{0c}$: ``on a typical
  (65--70$^\circ$F) day, ridership is around 5000, likely 3000--7000.''
  Nobody has intuition about rides at 0$^\circ$F. Centering makes the
  prior easy to elicit and interpret (and \texttt{rstanarm} uses this
  scale).
\end{frame}

\begin{frame}{Answer 9.3}
  \textbf{(a)} The credible interval reflects only \emph{posterior
  variability} (where the line is). The prediction interval adds
  \emph{sampling variability}: individual days scatter around the line
  with SD $\sigma \approx 1281$.
  \vspace{0.4em}

  \textbf{(b)} The \textbf{prediction interval}: the manager faces one
  specific day, not the long-run average.
  \vspace{0.4em}

  \textbf{(c)} The credible interval for the mean shrinks toward a point
  (the line becomes certain). The prediction interval barely changes: it
  stays near mean $\pm 1.96\,\sigma \approx \pm 2500$ rides, because
  day-to-day scatter never goes away. (With 500 days, posterior
  variability is already small next to $\sigma$.)
\end{frame}

\begin{frame}{Answer 9.4}
  \begin{itemize}
    \item A slope is learned from \emph{variation in $X$}. Thirty winter
      days cover a narrow, cold temperature range, so they say little
      about what happens on warm days.
    \item The information in data depends on \emph{what} was observed,
      not just \emph{how much}.
    \item Each phase updates the previous posterior. As data span more
      temperatures, the posterior for $\beta_1$ moves and narrows,
      ending concentrated near 80.
    \item Updating in phases gives the same final posterior as using all
      500 days at once.
  \end{itemize}
\end{frame}

\begin{frame}{Answer 9.5}
  \textbf{(a)} $100 \pm 2(40) = (20,\ 180)$: each extra degree adds
  roughly 20 to 180 rides. The prior is sure the relationship is positive
  but vague about its size.
  \vspace{0.4em}

  \textbf{(b)} $1/0.0008 = 1250$ rides: days typically scatter about
  1250 rides around the line.
  \vspace{0.4em}

  \textbf{(c)} $-2194 + 82.2 \times 75 = -2194 + 6165 = 3971$ rides
  (the book reports $\approx 3968$; the gap is rounding in the
  coefficients). \\
  A 10$^\circ$F difference: $82.2 \times 10 = 822$ more rides on the
  70$^\circ$F day. The intercept cancels.
\end{frame}

\begin{frame}{Answer 9.6}
  \textbf{(a)}
  \begin{center}
  \begin{tabular}{@{}ccc@{}}
    \toprule
    Draw & $\mu = \beta_0 + 75\beta_1$ & $Y' = \mu + \sigma z$ \\
    \midrule
    1 & $-2200 + 6150 = 3950$ & $3950 + 650 = 4600$ \\
    2 & $-2100 + 6075 = 3975$ & $3975 - 1250 = 2725$ \\
    3 & $-2300 + 6300 = 4000$ & $4000 + 256 = 4256$ \\
    \bottomrule
  \end{tabular}
  \end{center}
  \textbf{(b)} The $\mu$'s build the credible interval for the
  \emph{average}; the $Y'$'s build the prediction interval. The $Y'$'s
  are far more spread out (2725 to 4600 vs.\ 3950 to 4000) because they
  add day-to-day scatter $\sigma z$ on top of the uncertainty in the line.
  With 20{,}000 draws, this is exactly what \texttt{posterior\_predict()}
  does.
\end{frame}

%---------------- Chapter 15 answers
\begin{frame}{Answer 15.1}
  \begin{enumerate}
    \item \textbf{Grouped by school.} Students in one school share
      teachers, resources and neighborhood.
    \item \textbf{Grouped by patient.} Repeated measures on one person
      reflect that person's physiology.
    \item \textbf{Not grouped.} Independent random sample, one response
      each.
    \item \textbf{Grouped by city.} Each city has its own climate (and
      consecutive days are also correlated over time).
  \end{enumerate}
  \vspace{0.4em}
  The common thread: knowing one observation from a group tells you
  something about the others in that group.
\end{frame}
\begin{frame}{Answer 15.2}
  \begin{itemize}
    \item Runners differ a lot in baseline speed (between-runner
      variability). Complete pooling treats those differences as random
      noise, inflating $\sigma$ and blurring the age effect.
    \item Comparisons \emph{across} runners (a fast 60-year-old vs.\ a slow
      55-year-old) are mixed with changes \emph{within} each runner, which
      is what ``slowing with age'' actually means.
    \item Violated assumption: \textbf{independence}. Times from the same
      runner are correlated, and runner identity is ignored.
    \item Result: a misleading conclusion about the trend and its
      significance.
  \end{itemize}
\end{frame}

\begin{frame}{Answer 15.3}
  \textbf{(a)} Five races are very little information. The model
  extrapolates runner 1's short-term improvement for years, ignoring
  what the other 35 runners show: people in their 50s--60s tend to slow
  down. It uses no information from other groups.
  \vspace{0.6em}

  \textbf{(b)} It has a separate model for each sampled runner and no model
  of the \emph{population} of runners. There is nothing to say how a new
  runner's intercept and slope relate to the others, so it cannot
  generalize beyond the sample.
\end{frame}

\begin{frame}{Answer 15.4}
  \textbf{(a)} \emph{Within}: how much one runner's times vary year to year
  (small for runner 29, large for runner 17). \emph{Between}: how much
  runners differ from each other in typical speed (29 vs.\ 17).
  \vspace{0.4em}

  \textbf{(b)} Each runner's ability persists: a runner who was fast last
  year is likely fast this year. So knowing one time tells us about the
  next.
  \vspace{0.4em}

  \textbf{(c)} It models each runner's own trend (respects between-group
  differences) \emph{and} treats runners as drawn from one population, so
  each estimate is pulled toward the overall trend. ``Learn from one
  another while celebrating our individuality.''
\end{frame}

\begin{frame}{Answer 15.5}
  \textbf{(a)} Both runners slow down by \textbf{+1 min/year}.
  \vspace{0.4em}

  \textbf{(b)} $\bar x = 58.5$, $\bar y = 76$.
  \begin{itemize}
    \item $x - \bar x$: $-3.5, -2.5, -1.5, 1.5, 2.5, 3.5$
    \item $y - \bar y$: $4, 5, 6, -6, -5, -4$
    \item $\sum (x-\bar x)(y-\bar y) = -71$, \quad
      $\sum (x-\bar x)^2 = 41.5$
  \end{itemize}
  $b = -71/41.5 \approx -1.71$ min/year: runners appear to get
  \emph{faster} with age!
  \vspace{0.4em}

  \textbf{(c)} The pooled line compares a slow younger runner with a fast
  older runner: differences \emph{between} runners masquerade as an age
  effect. Within each runner, time goes up with age. Ignoring the grouping
  can hide or even reverse a trend (as in Q15.1).
\end{frame}

\begin{frame}{Answer 15.6}
  \textbf{(a)} $w = 0.5$: $0.5(80) + 0.5(90) = 85$ min. \\
  $w = 0.9$: $0.9(80) + 0.1(90) = 81$ min.
  \vspace{0.4em}

  \textbf{(b)} X (2 races) $\to w = 0.5$; Y (10 races) $\to w = 0.9$.
  With more data, a runner's own mean is more reliable, so it gets more
  weight. With little data, the estimate borrows more from the other
  runners (it is pulled more toward the overall mean).
  \vspace{0.4em}

  \textbf{(c)} $w = 0$: complete pooling (everyone gets 90). $w = 1$: no
  pooling (each runner keeps their own mean). Partial pooling is
  everything in between. Chapter 16 shows how $w$ also depends on the
  within- and between-runner variability.
\end{frame}

%---------------- Multiple choice answers
\begin{frame}{Multiple-choice answers}
  \small
  \textbf{MCQ 1: (C) $\hat R$.} It compares the variability across all
  chains combined with the variability within each chain.
  $N_{\text{eff}}/N$, autocorrelation and a single trace plot look at
  chains one at a time and can miss chains stuck in different places.
  \vspace{0.5em}

  \textbf{MCQ 2: (B) 0.25.} $\alpha = \min\{1,\ 0.1/0.4\} = 0.25$.
  Only the ratio matters; the normalizing constant cancels.
  \vspace{0.5em}

  \textbf{MCQ 3: (C).} The prediction interval combines posterior
  variability and sampling variability. It is \emph{wider} than the
  credible interval for the mean (A is false), and it never shrinks to a
  point because days scatter around the line with SD $\sigma$ (D is
  false).
  \vspace{0.5em}

  \textbf{MCQ 4: (C).} Complete pooling has one universal model;
  the hierarchical model describes the population of runners. The
  no-pooled model only knows the 36 sampled runners.
\end{frame}

\end{document}
